% chapters/requirements/requirements.tex
% Capítulo 3: Análisis de requisitos del backend del SCDEE/B
% ============================================================================

Este capítulo recoge los requisitos funcionales y no funcionales del subsistema
\textit{backend} del \ac{scdee}. Los requisitos se organizan en 16~subsistemas
funcionales (RF) y 15~requisitos no funcionales (RNF). Cada requisito funcional se
describe con su entrada (petición o evento que lo desencadena), proceso (lógica
que ejecuta el \textit{backend}) y salida (respuesta o efecto producido). La
\cref{tab:resumen_requisitos} presenta un resumen cuantitativo de los
subsistemas y el número de requisitos funcionales de cada uno. Los \acsp{rnf}
establecen restricciones transversales de rendimiento, seguridad, despliegue y
mantenibilidad.

El análisis parte de la descripción del proyecto proporcionada por el tutor, de las
extensiones funcionales definidas durante la fase de especificación y de las decisiones
técnicas tomadas para resolver las ambigüedades identificadas.

Al tratarse de un \ac{tfg} coordinado, este catálogo se centra exclusivamente en los
requisitos del backend. Los requisitos de presentación, interacción y experiencia de
usuario son responsabilidad del \ac{tfg} de frontend \cite{soler2026frontend}. Se puede
consultar el catálogo global de requisitos en el anexo \ref{ap:global_requirements}.
